%------------------------------------------------------------------------------------
%  SECTION 1.3 : GESTION DES INTERLIGNES ET MICROTYPOGRAPHIE
%------------------------------------------------------------------------------------
\section{Interlignage Contrôlé et Microtypographie Avancée}
\label{sec:interlignes_microtype}

La lisibilité d'un document long repose sur une harmonie subtile entre la taille de la police, la longueur de la ligne et l'espacement vertical entre deux lignes consécutives.
Trop serré, le texte fatigue l'œil lors des retours à la ligne ; trop espacé, la page semble déstructurée et les paragraphes perdent leur cohésion.

\subsection{Le défi classique de l'interligne 1,5 et la réponse de la classe}
\label{subsec:defi_interligne}

La majorité des universités et jurys de thèse exigent un interligne de 1,5 pour le confort de lecture et l'annotation manuscrite des épreuves.
Le paquet standard \textsf{setspace} fournit la commande \texttt{\textbackslash onehalfspacing}.

Cependant, en \LaTeX{} standard (\emph{vanilla}), l'activation globale de cet interligne engendre un effet visuel désastreux :
lorsqu'un titre de section ou de chapitre s'étend sur deux lignes, l'interligne 1,5 lui est également appliqué !
Le titre apparaît alors artificiellement distendu et déconnecté de son texte.

\begin{boiteSynthese}{Protection automatique des titres dans ClasseGuillaume}
  Pour pallier cette anomalie, \texttt{ClasseGuillaume} intègre une règle typographique stricte au sein des polices KOMA-Script :
  \begin{verbatim}
  \addtokomafont{section}{\setstretch{1}\selectfont}
  \addtokomafont{subsection}{\setstretch{1}\selectfont}
  \addtokomafont{subsubsection}{\setstretch{1}\selectfont}
  \end{verbatim}
  Ainsi, le corps du mémoire bénéficie de la fluidité de l'interligne 1,5, tandis que l'ensemble des titres longs conservent un interligne simple compact, rigoureux et professionnel.
\end{boiteSynthese}

\subsection{Microtypographie et élimination des lignes orphelines et veuves}
\label{subsec:microtypographie}

Un document scientifique soigné ne doit jamais laisser une ligne isolée au bas d'une page (ligne orpheline) ou en haut de la page suivante (ligne veuve).
De même, une césure coupant un mot à la toute dernière ligne d'une page est considérée comme un manquement aux règles de l'art typographique.

Le paquet \textsf{microtype} associé à la configuration interne de \texttt{ClasseGuillaume} résout ces écueils par un ajustement microscopique de l'espacement inter-mots et par des pénalités strictes :

\begin{itemize}
  \item \texttt{\textbackslash clubpenalty = 10000} : Interdit formellement la présence d'une ligne orpheline en bas de page.
  \item \texttt{\textbackslash widowpenalty = 10000} : Interdit le rejet d'une ligne veuve en haut de page.
  \item \texttt{\textbackslash displaywidowpenalty = 10000} : Protège l'agencement des formules mathématiques détachées.
  \item \texttt{\textbackslash brokenpenalty = 10000} : Empêche la césure d'un mot à cheval sur deux pages consécutives.
  \item \texttt{\textbackslash interfootnotelinepenalty = 9999} : Empêche qu'une note de bas de page ne soit scindée arbitrairement sur deux pages.
\end{itemize}

Grâce à ces paramètres, le compilateur Lua\LaTeX{} recalcule subtilement les blancs typographiques de chaque paragraphe pour produire des pages équilibrées sans rupture visuelle.

